\subsection{合成实验}

我们首先在若干合成文本数据集上开展实验，这些数据集由五个作用于二进制串的 Tomita 文法，以及一个定义在三值字母表 $\Sigma_{M} = \{\mathtt{\ (\ ,\ \starchar\ ,\ )\ }\}$ 上的上下文无关的 ``Motzkin'' 文法组成~\citep{tomita1982, alexander2018}。后者由所有括号正确配对的串组成，且对 $\mathtt{\starchar}$ 字符不加任何约束。

在每种情形下，我们都在来自该文法、长度受限的串上训练 u-MPS，然后从训练好的 u-MPS 中采样未见长度的新串，并以采样串中匹配该文法的比例来评估模型。采样有两种形式：其一是固定长度 $n$ 的采样（对应 $R = \Sigma^n$）；其二是字符补全采样（character completion sampling），即将参考串中的一个字符遮蔽，用前缀和后缀 $p$ 和 $s$ 来猜测它（对应 $R = p \Sigma s$）。虽然很容易设想更一般的采样实验，但我们之所以选择这些任务，是因为它们能够与多种基线进行直接比较，包括（单向和双向的）LSTM、HMM 以及 Transformer。

虽然 u-MPS 可以通过算法~\ref{alg:sampling} 轻松实现无偏的固定长度采样，但我们发现单向 LSTM 基线需要在其输入中加入额外的位置编码（positional encoding），以避免在采样超出训练中所见最长长度时输出文本迅速退化。在采样时，我们根据期望的采样长度调整与该编码相关的长度尺度，使得采样的最后一步总是对应同一个位置编码向量。

我们使用负对数似然（NLL）损失和 Adam~\citep{kingma2015} 优化器，通过梯度下降训练 u-MPS、LSTM、HMM 以及小型 Transformer 模型。在每个实验中，我们使用隐单元数为 $D=20$ 和 $D=50$ 的模型，其中 LSTM 取单层，Transformer 取两层、4 个注意力头。每个 $D$ 值都进行五次独立试验，并用最终验证损失选出用于生成样本的最佳模型。我们在 $10^{-2}$ 与 $10^{-4}$ 之间使用分段常数学习率，并使用早停（early stopping）来确定训练结束的时机。

在 Tomita 实验（表~\ref{tab:tomita}）中，u-MPS 表现出色，在许多情形下对语言内未见长度的串采样达到了完美的精度。这不仅对较简单的 Tomita 3 和 Tomita 4 文法成立，对更难的 Tomita 5 也成立，其中合法串满足包含偶数个 \texttt{0} 和偶数个 \texttt{1} 的非局域约束。HMM 与 LSTM 也取得了相当高的采样精度，但二者随序列长度增加而退化的速度快于 u-MPS，而 Transformer 的表现最差。

在上下文无关的 Motzkin 语言（表~\ref{tab:motzkin}）上也看到了类似结果：与 Tomita 实验类似的固定长度采样任务与一个字符补全任务相配对。后一任务必须使用单独的双向 LSTM，因为单向 LSTM 无法利用未来上下文信息。相比之下，训练好的 u-MPS 模型可以在这两种设定中直接使用而无需任何任务特定的 adaptation，也可用于更一般的、涉及相连或不相连缺失文本区域的句子补全任务（标准 RNN 模型难以处理这些任务）。u-MPS 在泛化和再现 Motzkin 串的结构方面显著优于单向 LSTM、HMM 与 Transformer，能够以超过 90\% 的精度采样出长度达到训练所见 3 倍以上的串。只有双向 LSTM 在较小训练集的字符补全实验中胜过了 u-MPS。

鉴于 HMM 能够精确再现与 Tomita 语言以及（有界长度的）Motzkin 语言相关联的分布，u-MPS 在实践中仍能更轻易地学到这类分布，这一点颇为令人惊讶。同样有些令人惊讶的是 Transformer 模型的糟糕表现，其原因很可能在于所用字符串数据集的规模较小。

我们还对在 CPU 或 GPU 上运行的 u-MPS 的顺序求值与并行求值的相对运行时间进行了基准测试（表~\ref{tab:runtime}）。对于具有 $D$ 个隐状态的 u-MPS 和长度为 $n$ 的串，RNN 式的顺序求值具有 $\mathcal{O}(n)$ 的并行深度和 $\mathcal{O}(nD^2)$ 的代价，而并行求值则以 $\mathcal{O}(\log n)$ 的并行深度和 $\mathcal{O}(nD^3)$ 的代价取代之。这暗示在存在硬件加速时并行求值对总运行时间有好处，表~\ref{tab:runtime} 中的运行时间证实了这一点。

